% Part: second-order-logic
% Chapter: metatheory
% Section: loewenheim-skolem

\documentclass[../../../include/open-logic-section]{subfiles}

\begin{document}

\olfileid{sol}{met}{lst}

\olsection{شکست قضیهٔ لوونهایم--اسکولم در منطق مرتبهٔ دوم}

\begin{explain}
قضیهٔ (نزولیِ) لوونهایم--اسکولم می‌گوید هر مجموعه از !!{sentence}s که
مدلی نامتناهی داشته باشد، مدلی~!!a{enumerable} دارد. این قضیه نیز نتیجهٔ
قضیهٔ تمامیت است: اثبات تمامیت برای هر مجموعهٔ سازگار از !!{sentence}s
مدلی می‌سازد و آن مدل !!{enumerable} است. قضیهٔ صعودیِ
لوونهایم--اسکولم نیز وجود دارد که تضمین می‌کند اگر مجموعه‌ای از
!!{sentence}s مدلی~!!a{denumerable} داشته باشد، مدلی~!!a{nonenumerable}
هم دارد. هر دو قضیه در منطق مرتبهٔ دوم شکست می‌خورند.
\end{explain}


\begin{thm}
\ollabel{thm:sol-no-ls} قضیهٔ لوونهایم--اسکولم در منطق مرتبهٔ دوم شکست
می‌خورد: !!{sentence}sی وجود دارند که مدل نامتناهی دارند اما هیچ مدل
!!{enumerable}ی ندارند.
\end{thm}

\begin{proof}
به یاد آورید که
\[
\fn{Count} \ident \lexists[z][\lexists[u][\lforall[X][((X(z) \land
      \lforall[x][(X(x) \lif X(u(x)))]) \lif \lforall[x][X(x)])]]]
\]
در !!a{structure}ی چون~$\Struct{M}$ صادق است اگر و تنها اگر
$\Domain{M}$ !!{enumerable} باشد؛ پس~$\lnot\fn{Count}$ در~$\Struct{M}$
صادق است اگر و تنها اگر~$\Domain{M}$ !!{nonenumerable} باشد. چنین
!!{structure}sی وجود دارند: هر مجموعهٔ !!{nonenumerable}ی را، برای
نمونه~$\Pow{\Nat}$ یا~$\Real$، به‌عنوان !!{domain} بگیرید. پس
$\lnot \fn{Count}$ مدل‌های نامتناهی دارد اما هیچ مدل !!{enumerable}ی
ندارد.
\end{proof}

\begin{thm}
!!{sentence}sی وجود دارند که !!{denumerable} مدل دارند اما هیچ مدل
!!{nonenumerable}ی ندارند.
\end{thm}

\begin{proof}
$\fn{Count} \land \fn{Inf}$ در~$\Nat$ صادق است اما در هیچ
!!{structure}~$\Struct{M}$ با !!{nonenumerable}بودنِ~$\Domain{M}$
صادق نیست.
\end{proof}


\end{document}
